\chapter{Background}

\section{Retrieval-Augmented Generation}

Retrieval-Augmented Generation was introduced by \textcite{rag2020} as a method for grounding language model outputs in external document collections. A retriever---typically a dense bi-encoder---embeds queries and documents into a shared vector space and retrieves the most similar documents by nearest-neighbour search. A generator then conditions on both the query and the retrieved documents to produce the final answer. The key advantage over purely parametric approaches is that the knowledge base can be updated by re-indexing documents without retraining the model, and retrieved passages provide an auditable evidence trail for generated claims.

The practical significance of RAG for domain-specific question answering is substantial. A general-purpose LLM trained on internet-scale data has no reliable knowledge of institution-specific details---current staff contacts, procedural requirements, room assignments---that change frequently and were not prominently represented in training corpora. RAG provides a principled mechanism to inject this dynamic, local knowledge at inference time.

\section{Self-RAG}

A notable limitation of standard RAG is that the system blindly generates a response regardless of whether the retrieved passages are actually relevant or sufficient to answer the question. \textcite{selfrag2024} address this with Self-RAG, a framework where the model generates special \emph{reflection tokens} to decide when retrieval is needed, assess the relevance of each retrieved passage, and evaluate whether its own output is supported by the evidence. In the original formulation, the model is fine-tuned end-to-end to produce these tokens as part of the standard generation process.

Our implementation adopts the spirit of Self-RAG without task-specific fine-tuning. A general-purpose Qwen3 model is prompted to perform each reflection step as a structured classification task, producing JSON-formatted judgements that control the pipeline's branching logic. This trades the elegance of the original single-model approach for deployment simplicity: no specialised training data or fine-tuning infrastructure is required.

\section{Hybrid Retrieval}

Dense retrieval using bi-encoder models captures semantic similarity: query and document are encoded independently and compared by inner product or cosine similarity. This works well when the query and relevant passage use different vocabulary for the same concept. However, dense models can fail on exact-match queries---specific names, course codes, or building identifiers---if these terms are not well-represented in the embedding space.

Sparse retrieval methods such as BM25 \parencite{bm25Robertson2009} rely on lexical overlap and inverse document frequency weighting, excelling precisely where dense retrieval struggles. The complementary nature of these two signals makes their combination consistently stronger than either alone. Reciprocal Rank Fusion (RRF) \parencite{rrf2009} provides a simple, parameter-free method for merging two ranked lists: results are re-ranked by summing reciprocal rank scores, requiring no normalisation of heterogeneous similarity scores.

Once a candidate set is assembled through retrieval, a cross-encoder reranker \parencite{sentencetransformers2019} jointly encodes the query and each candidate passage, enabling full attention-based interaction modelling. Cross-encoders are substantially more accurate than bi-encoders at judging fine-grained relevance, but are too expensive to apply to the entire corpus; they are therefore used only to re-rank a small pre-filtered candidate set.
